\documentclass[11pt,fontset=fandol]{ctexart}
\usepackage[a4paper,margin=23mm]{geometry}
\usepackage{amsmath,amssymb,mathtools,bm,mathrsfs,esint}
\usepackage{graphicx,longtable,booktabs,array,calc}
\usepackage[normalem]{ulem}
\usepackage{xcolor,hyperref}
\hypersetup{hidelinks}
\setcounter{secnumdepth}{-1}
\newcounter{none}
\providecommand{\tightlist}{\setlength{\itemsep}{0pt}\setlength{\parskip}{0pt}}
\providecommand{\passthrough}[1]{#1}
\usepackage{amssymb,mathrsfs,esint,bm,newunicodechar}
\allowdisplaybreaks
\setlength{\emergencystretch}{3em}
\providecommand{\llangle}{\langle\!\langle}
\providecommand{\rrangle}{\rangle\!\rangle}
\newunicodechar{∞}{\ensuremath{\infty}}
\newunicodechar{∈}{\ensuremath{\in}}
\newunicodechar{∑}{\ensuremath{\sum}}
\newunicodechar{∫}{\ensuremath{\int}}
\newunicodechar{≤}{\ensuremath{\leq}}
\newunicodechar{≥}{\ensuremath{\geq}}
\newunicodechar{≠}{\ensuremath{\neq}}
\newunicodechar{→}{\ensuremath{\to}}
\newunicodechar{⇒}{\ensuremath{\Rightarrow}}
\newunicodechar{α}{\ensuremath{\alpha}}
\newunicodechar{β}{\ensuremath{\beta}}
\newunicodechar{γ}{\ensuremath{\gamma}}
\newunicodechar{λ}{\ensuremath{\lambda}}
\newunicodechar{φ}{\ensuremath{\varphi}}
\newunicodechar{ε}{\ensuremath{\varepsilon}}
\newunicodechar{σ}{\ensuremath{\sigma}}
\newunicodechar{θ}{\ensuremath{\theta}}
\newunicodechar{δ}{\ensuremath{\delta}}
\newunicodechar{η}{\ensuremath{\eta}}
\newunicodechar{π}{\ensuremath{\pi}}
\newunicodechar{ξ}{\ensuremath{\xi}}
\newunicodechar{Σ}{\ensuremath{\Sigma}}
\newunicodechar{Δ}{\ensuremath{\Delta}}
\newunicodechar{∂}{\ensuremath{\partial}}
\newunicodechar{∀}{\ensuremath{\forall}}
\newunicodechar{∃}{\ensuremath{\exists}}
\newunicodechar{√}{\ensuremath{\surd}}
\newunicodechar{−}{\ensuremath{-}}
\newunicodechar{μ}{\ensuremath{\mu}}
\newunicodechar{ρ}{\ensuremath{\rho}}
\newunicodechar{τ}{\ensuremath{\tau}}
\newunicodechar{ω}{\ensuremath{\omega}}
\newunicodechar{Ω}{\ensuremath{\Omega}}
\newunicodechar{χ}{\ensuremath{\chi}}
\newunicodechar{ψ}{\ensuremath{\psi}}
\newunicodechar{Γ}{\ensuremath{\Gamma}}
\newunicodechar{Φ}{\ensuremath{\Phi}}
\newunicodechar{Λ}{\ensuremath{\Lambda}}
\newunicodechar{ℝ}{\ensuremath{\mathbb{R}}}
\newunicodechar{ℤ}{\ensuremath{\mathbb{Z}}}
\newunicodechar{ℕ}{\ensuremath{\mathbb{N}}}
\newunicodechar{△}{\ensuremath{\triangle}}
\newunicodechar{⊗}{\ensuremath{\otimes}}
\begin{document}
\section{2024年杭州电子科技大学601数学分析真题及解析(4页版)}\label{ux5e74ux676dux5ddeux7535ux5b50ux79d1ux6280ux5927ux5b66601ux6570ux5b66ux5206ux6790ux771fux9898ux53caux89e3ux67904ux9875ux7248}

来源：\texttt{HDU/601数学分析/2024/2024年杭州电子科技大学601数学分析真题及解析(4页版).pdf}。原
PDF 页序：3, 4。

\section{杭州电子科技大学2024年数学分析(回忆版)真题解析}\label{ux676dux5ddeux7535ux5b50ux79d1ux6280ux5927ux5b662024ux5e74ux6570ux5b66ux5206ux6790ux56deux5fc6ux7248ux771fux9898ux89e3ux6790}

\subsection{一.填空题(8题，每题5分，共40分)}\label{ux4e00.ux586bux7a7aux98988ux9898ux6bcfux98985ux5206ux517140ux5206}

\begin{enumerate}
\def\labelenumi{\arabic{enumi}.}
\tightlist
\item
  \(\lim _ { n \to \infty } ( n ^ { 2 0 2 4 } + 2 ^ { n } ) ^ { \frac { 1 } { n } } =\)
\end{enumerate}

\[
\lim _ { n \to \infty } ( n ^ { 2 0 2 4 } + 2 ^ { n } ) ^ { \frac { 1 } { n } } = \lim _ { n \to \infty } \Biggl [ 2 ^ { n } ( { \frac { n ^ { 2 0 2 4 } } { 2 ^ { n } } } + 1 ) \Biggr ] ^ { \frac { 1 } { n } } = 2 \lim _ { n \to \infty } ( { \frac { n ^ { 2 0 2 4 } } { 2 ^ { n } } } + 1 ) ^ { \frac { 1 } { n } } = 2 ,
\]

这里，
\(\lim _ { n \to \infty } ( { \frac { n ^ { 2 0 2 4 } } { 2 ^ { n } } } + 1 ) = 1\)
\(\lim _ { n \to \infty } { \frac { 1 } { n } } = 0\) ，故
\(\lim _ { n \to \infty } ( { \frac { n ^ { 2 0 2 4 } } { 2 ^ { n } } } + 1 ) ^ { \frac { 1 } { n } } = 1 ^ { 0 } = 1\)

\begin{enumerate}
\def\labelenumi{\arabic{enumi}.}
\setcounter{enumi}{1}
\tightlist
\item
  \(\int \frac { x ^ { 3 } } { \sqrt { 1 - x ^ { 2 } } } d x =\)
\end{enumerate}

\[
\int \frac{x^3}{\sqrt{1-x^2}}\,dx
= -\int (1-u^2)\,du
= \frac{u^3}{3}-u+C
= \frac{(1-x^2)^{3/2}}{3}-\sqrt{1-x^2}+C,
\qquad u=\sqrt{1-x^2}.
\]

\begin{enumerate}
\def\labelenumi{\arabic{enumi}.}
\setcounter{enumi}{2}
\tightlist
\item
  \(y = 2 ^ { x }\) 的麦克劳林级数的 \(x ^ { 2 0 2 4 }\) 的系数为
\end{enumerate}

\(y = 2 ^ { x }\) 的麦克劳林级数的 \(x ^ { 2 0 2 4 }\) 的系数为
\({ \frac { y ^ { ( 2 0 2 4 ) } ( 0 ) } { 2 0 2 4 ! } } = { \frac { ( \ln 2 ) ^ { 2 0 2 4 } } { 2 0 2 4 ! } }\)
4.
\(\lim _ { n \to \infty } { \frac { ( n + 1 ) \sin x } { x n + 4 } }\)
的间断点为

\[
\lim _ { n \to \infty } { \frac { ( n + 1 ) \sin x } { x n + 4 } } = \left\{ { \begin{array} { l } { 0 , x = 0 } \\ { \displaystyle { \frac { \sin x } { x } } , x \neq 0 } \end{array} } , \right.
\]

其间断点为 \(x = 0\) ·

5.求 \(x ^ { 2 } + y ^ { 2 } + z ^ { 2 } = 6\) 与
\(x ^ { 2 } + y ^ { 2 } = z\) 的交线的切线方程.

联立方程 \(x^2+y^2+z^2=6\) 与 \(x^2+y^2=z\)，得到 \(z^2+z-6=0\)。由于
\(z=x^2+y^2\geq0\)，所以交线满足 \(z=2\)、\(x^2+y^2=2\)。可取参数方程
\(x=\sqrt2\cos t\)、\(y=\sqrt2\sin t\)、\(z=2\)。在交线上任一点
\((x_0,y_0,2)\)，切向量为 \((-y_0,x_0,0)\)，切线方程可写为

\[
\begin{cases}
x_0x+y_0y=2,\\
z=2.
\end{cases}
\]

本段由助手依据原方程及可识别的参数解法整理，消除了识别稿中损坏的括号。

\subsection{6.没回忆到.}\label{ux6ca1ux56deux5fc6ux5230.}

\begin{enumerate}
\def\labelenumi{\arabic{enumi}.}
\setcounter{enumi}{6}
\tightlist
\item
  \(f ( x , y , z ) = x ^ { 2 } + x y + \cdots\) ，从
  \(P _ { 0 } ( - 1 , 0 , 2 )\) 到P(2,1,1)的方向导数.(没回忆全) 8.L为
  \(x ^ { 2 } + y ^ { 2 } = 1\) 上半圆周，则
  \(\int _ { L } ( x ^ { 2 } + y ^ { 2 } ) d s =\) 的
\end{enumerate}

L为 \(x ^ { 2 } + y ^ { 2 } = 1\) 上半圆周，故
\(\int _ { L } ( x ^ { 2 } + y ^ { 2 } ) d s = \int _ { L } d s = \pi\)

\subsection{二.计算题(共25分)}\label{ux4e8c.ux8ba1ux7b97ux9898ux517125ux5206}

1.(5分）
\(\lim _ { x \to 0 } { \frac { x - \sin { x } } { x ( e ^ { x } - 1 ) } }\)

\[
\lim _ { x \to 0 } { \frac { x - \sin x } { x ( e ^ { x } - 1 ) } } = \lim _ { x \to 0 } { \frac { \frac { 1 } { 6 } x ^ { 3 } } { x ^ { 2 } } } = 0 .
\]
\end{document}
